\documentclass{fiche}
\metadonnees{10}{Camper sous la pluie}{Bloc 2 — Ce qui est possible : futur et modaux}

\begin{document}
\entete

\objectifs{%
\textbf{Langue} : bilan des modaux ; modal \emph{+ have} + participe ; expression écrite.\par
\textbf{Texte} : Jerome K. Jerome, \emph{Three Men in a Boat} (1889), chapitre \textsc{ii}.\par
\textbf{Durée} : 1\,h\,30 — exercices 35 min, texte et questions 30 min, version 25 min.}

%% ===================================================================
\exercice[13 min]{Bilan des modaux}

\rappel{\textbf{Rappel.} Un modal dit ce que le locuteur pense de l'action : possible,
obligatoire, souhaitable, probable. Il ne change jamais de forme et se construit sans
\emph{to}.}

\consigne{Complétez avec le modal qui convient : \emph{can}, \emph{could}, \emph{may},
\emph{might}, \emph{must}, \emph{mustn't}, \emph{needn't}, \emph{should}, \emph{had better}.}

\begin{anglais}
\begin{enumerate}[label=\arabic*.,itemsep=2.2mm,leftmargin=*]
  \item You \trou{can} never rouse Harris; there is no poetry about him.
  \item If your eyes are full of tears, you \trou{must} have been eating raw onions.
  \item We \trou{had better} pitch the tent before the light goes.
  \item \trou{May} I leave my coat in the boat?
  \item It \trou{might} rain tonight, but nobody is sure.
  \item You \trou{needn't} take the heavy stove; there is an inn at Chertsey.
  \item You \trou{mustn't} light a fire so near the tent.
  \item \trou{Could} you hold this rope for a moment?
  \item They \trou{should} arrive before dusk if the wind drops.
  \item That \trou{can't} be Bill: Bill is under the tent.
  \item In dry weather a tent \trou{can} be fixed in ten minutes.
  \item We \trou{must} be at Kingston on Saturday morning; the boat is booked.
  \item You \trou{may} sleep at an inn if you prefer; nobody is forcing you.
  \item He \trou{should} have listened to the warning about the weather.
\end{enumerate}
\end{anglais}

\note[Les quatre emplois]{Capacité ou possibilité générale : 1, 8, 11. Permission : 4, 13.
Obligation ou conseil : 3, 6, 7, 12, 14. Degré de certitude : 2, 5, 9, 10. Un même modal
change de valeur selon le contexte : \emph{can} est tantôt la capacité (11), tantôt la
demande polie (8) ; \emph{may} tantôt la permission (13), tantôt la possibilité.}

%% ===================================================================
\exercice[10 min]{Modal \emph{+ have} + participe}

\rappel{\textbf{Rappel.} Le modal reste au présent ; c'est l'infinitif qui passe au parfait.
\emph{Must have} : déduction. \emph{Should have} : regret. \emph{Could have} : possibilité
non réalisée.}

\consigne{Traduisez la partie entre crochets.}

\begin{anglais}
\begin{enumerate}[label=\arabic*.,itemsep=3mm,leftmargin=*]
  \item {}[Il aurait dû] fix the tent in the open.
    \lignes{1}
    \corr{He should have fixed \emph{ou} He ought to have fixed}
  \item {}[Ils ont sûrement oublié] the mallet.
    \lignes{1}
    \corr{They must have forgotten}
  \item {}[Elle n'a pas pu entendre] us: she was asleep.
    \lignes{1}
    \corr{She can't have heard}
  \item {}[J'aurais pu t'aider] if you had asked.
    \lignes{1}
    \corr{I could have helped you}
  \item {}[Tu n'aurais pas dû] pull the ropes so hard.
    \lignes{1}
    \corr{You should not have pulled \emph{ou} You ought not to have pulled}
  \item {}[Il a peut-être perdu] his way in the dark.
    \lignes{1}
    \corr{He may have lost \emph{ou} He might have lost}
  \item {}[Ce n'était pas la peine de venir] : the work was done.
    \lignes{1}
    \corr{You need not have come}
  \item {}[Nous aurions dû écouter] the warning.
    \lignes{1}
    \corr{We should have listened to}
\end{enumerate}
\end{anglais}

\note[Une forme rare et utile]{\emph{Need not have} + participe (7) dit qu'on a fait une
chose inutile : \og ce n'était pas la peine de venir --- et pourtant tu es venu \fg. À
distinguer de \emph{did not need to come}, qui laisse entendre qu'on n'est pas venu.}

%% ===================================================================
\newpage
\section{Texte}

\noindent\small\textit{Trois amis --- le narrateur, George et Harris --- préparent une
remontée de la Tamise en barque. Reste à trancher : coucheront-ils sous la tente ou à
l'auberge ? Le narrateur vient de faire l'éloge lyrique du campement au clair de lune.
Harris a répondu : \og Et quand il pleut ? \fg}\normalsize

\begin{texte}
You can never \gl[to rouse]{rouse}{éveiller, émouvoir} Harris. There is no poetry about
Harris --- no wild \gl[a yearning]{yearning}{une aspiration ardente} for the unattainable.
Harris never ``weeps, he knows not why.'' If Harris's eyes fill with tears, you can bet it is
because Harris has been eating raw onions, or has put too much Worcester over his
\gl[a chop]{chop}{une côtelette}.

If you were to stand at night by the sea-shore with Harris, and say:

``Hark! do you not hear? Is it but the mermaids singing deep below the waving waters; or sad
spirits, chanting \gl[a dirge]{dirges}{un chant funèbre} for white corpses, held by
seaweed?'' Harris would take you by the arm, and say:

``I know what it is, old man; you've got a \gl[a chill]{chill}{un coup de froid}. Now, you
come along with me. I know a place round the corner here, where you can get a drop of the
finest Scotch whisky you ever tasted --- put you right in less than no time.''

Harris always does know a place round the corner where you can get something brilliant in the
drinking line. I believe that if you met Harris up in Paradise (supposing such a thing
likely), he would immediately greet you with:

``So glad you've come, old fellow; I've found a nice place round the corner here, where you
can get some really first-class \gl[nectar]{nectar}{le nectar}.''

In the present instance, however, as regarded the camping out, his practical view of the
matter came as a very timely hint. Camping out in rainy weather is not pleasant.

It is evening. You are wet through, and there is a good two inches of water in the boat, and
all the things are damp. You find a place on the banks that is not quite so
\gl[puddly]{puddly}{plein de flaques} as other places you have seen, and you land and
\gl[to lug out]{lug out}{traîner dehors} the tent, and two of you proceed to fix it.

It is \gl[soaked]{soaked}{trempé} and heavy, and it flops about, and tumbles down on you, and
clings round your head and makes you mad. The rain is pouring steadily down all the time. It
is difficult enough to fix a tent in dry weather: in wet, the task becomes
\gl[herculean]{herculean}{herculéen}. Instead of helping you, it seems to you that the other
man is simply playing the fool. Just as you get your side beautifully fixed, he gives it a
\gl[a hoist]{hoist}{une traction vers le haut} from his end, and spoils it all.
\end{texte}
\source{Jerome K. Jerome, \emph{Three Men in a Boat}, Bristol, Arrowsmith, 1889,
ch.~\textsc{ii}.}

\partiel[15 min]{Questions}
\consigne{\en{Answer in English, in full sentences.}}

\begin{enumerate}[label=\arabic*.,itemsep=2mm,leftmargin=*]
  \item \en{What kind of man is Harris, according to the narrator?}
    \lignes{2}
    \corr{A completely unpoetical one. He has no longing for the unattainable, he never weeps
    without a physical cause, and he always knows a place round the corner where one can get
    a drink.}
  \item \en{The narrator imagines meeting Harris in Paradise. What is the joke?}
    \lignes{2}
    \corr{Even in Paradise Harris would greet him with the address of a good bar. Nothing,
    not even eternity, can change him --- and the parenthesis \emph{supposing such a thing
    likely} suggests the narrator doubts Harris will get there at all.}
  \item \en{Why does the narrator call Harris's question a ``very timely hint''?}
    \lignes{2}
    \corr{Because it brings the discussion back to earth just in time. The narrator had been
    describing moonlit evenings; Harris reminds them that it also rains.}
  \item \en{Notice the change of tense at ``It is evening.'' What does the narrator do?}
    \lignes{3}
    \corr{He stops telling his own story and switches to the present with an indefinite
    \emph{you}: the scene becomes a general truth, something that happens to anybody who
    camps in the rain. The reader is put inside it.}
  \item \en{List the things that go wrong while the tent is being fixed.}
    \lignes{3}
    \corr{Everything is wet through and the boat has two inches of water; the ground is full
    of puddles; the tent is soaked, heavy, flops about, falls on them and clings round their
    heads; the rain never stops; and each man believes the other is playing the fool.}
  \item \en{Where does the comedy come from in the last paragraph?}
    \lignes{2}
    \corr{From accumulation and from the exaggerated contrast between the epic word
    \emph{herculean} and a task as small as putting up a tent.}
\end{enumerate}

\note[Le \emph{you} indéfini]{\emph{You are wet through}, \emph{you find a place},
\emph{makes you mad} : ce \emph{you} ne désigne pas le lecteur mais n'importe qui, comme le
\og on \fg{} français. Même valeur que \emph{one}, mais familier. Le repérer évite le
contresens en version.}

%% ===================================================================
\newpage
\section{Version}
\consigne{Traduisez le passage en français. Entre-temps la tente est tombée deux fois, le
souper s'est révélé être de l'eau de pluie, et les trois hommes se sont couchés.}

\begin{version}
There you dream that an elephant has suddenly sat down on your chest, and that the volcano
has exploded and thrown you down to the bottom of the sea --- the elephant still sleeping
peacefully on your \gl[a bosom]{bosom}{la poitrine}. You wake up and grasp the idea that
something terrible really has happened. Your first impression is that the end of the world
has come; and then you think that this cannot be, and that it is thieves and murderers, or
else fire, and this opinion you express in the usual method. No help comes, however, and all
you know is that thousands of people are kicking you, and you are being
\gl[to smother]{smothered}{étouffer}.

Somebody else seems in trouble, too. You can hear his faint cries coming from underneath your
bed. Determining, at all events, to sell your life dearly, you struggle frantically, hitting
out right and left with arms and legs, and \gl[lustily]{lustily}{à pleins poumons} yelling the
while, and at last something gives way, and you find your head in the fresh air. Two feet off,
you dimly observe a half-dressed \gl[a ruffian]{ruffian}{une brute} waiting to kill you, and
you are preparing for a life-and-death struggle with him, when it begins to
\gl[to dawn upon]{dawn upon}{apparaître peu à peu à} you that it's Jim.
\end{version}
\source{\emph{Ibid.}}

\lignes{22}

\corr{\textbf{Proposition de traduction.} Là, vous rêvez qu'un éléphant vient de s'asseoir
sur votre poitrine et que le volcan a fait éruption et vous a précipité au fond de la mer ---
l'éléphant dormant toujours paisiblement sur votre torse. Vous vous réveillez et saisissez
l'idée qu'il s'est réellement produit quelque chose d'épouvantable. Votre première impression
est que la fin du monde est arrivée ; puis vous vous dites que ce n'est pas possible, que ce
sont des voleurs et des assassins, ou bien le feu, et vous exprimez cette opinion par le
moyen ordinaire. Nul secours ne vient, pourtant, et tout ce que vous savez, c'est que des
milliers de gens vous donnent des coups de pied et qu'on est en train de vous étouffer.

Quelqu'un d'autre paraît en difficulté, lui aussi. Vous entendez ses faibles cris monter de
sous votre lit. Résolu, en tout état de cause, à vendre chèrement votre vie, vous vous
débattez frénétiquement, frappant à droite et à gauche des bras et des jambes, en hurlant à
pleins poumons, et enfin quelque chose cède, et vous retrouvez la tête au grand air. À deux
pieds de là, vous distinguez confusément une brute à demi vêtue qui attend de vous tuer, et
vous vous préparez à une lutte à mort avec elle, quand peu à peu il vous apparaît que c'est
Jim.}

\note[Choix de traduction 1]{\emph{you are being smothered} : passif progressif. Le français
n'en dispose pas et passe par \og on \fg{} + actif (\og on est en train de vous
étouffer \fg). C'est la solution la plus fréquente pour tout passif anglais sans complément
d'agent.}

\note[Choix de traduction 2]{\emph{this opinion you express in the usual method} : l'objet
est placé en tête pour l'emphase. Le français ne le peut pas sans lourdeur ; il rétablit
l'ordre normal, et compense par \og cette opinion \fg{} en position détachée si l'on veut
garder l'effet.}

\note[Choix de traduction 3]{\emph{Determining, at all events, to sell your life dearly} :
participe présent en tête de phrase, valeur de cause ou de circonstance. Le français préfère
souvent un participe passé ou un adjectif (\og Résolu à\ldots \fg) au participe présent, qui
sonne administratif.}

\note[Choix de traduction 4]{\emph{it begins to dawn upon you that} : \emph{to dawn} = poindre,
se lever (le jour). L'image est celle d'une aube : \og il commence à vous apparaître \fg,
\og la lumière se fait \fg. Ne pas chercher un verbe unique en français : l'étoffement est
obligatoire.}

\note[Choix de traduction 5]{\emph{yelling the while} : \emph{the while} = \og pendant ce
temps \fg, tour figé et un peu archaïque. Le français le fond dans le gérondif \og en
hurlant \fg.}

%% ===================================================================
\partiel[15 min]{Expression écrite}
\consigne{Vous êtes le troisième homme, celui qui écopait le bateau pendant que les deux
autres se disputaient. Rédigez la page de votre journal pour cette soirée : 120 mots. Employez
au moins quatre modaux différents, dont un \emph{should have} + participe.}
\lignes{12}

\corr{\textbf{Proposition.} \emph{Thursday, near Datchet.} We ought to have slept at an inn.
I said so this morning, and nobody listened. It began to rain at four and it has not stopped
since. While I was baling out the boat --- three inches of water, and my sleeve full of it
--- the other two were fighting with the tent. They could have fixed it in ten minutes in dry
weather; tonight they walked round it for a quarter of an hour, swearing at each other, until
it came down on their heads. Supper was rainwater with a little bread in it. My tobacco is
damp, so I cannot even smoke. Tomorrow we shall find a hotel. We must, or I shall go home by
train.}

\note[Barème]{Quatre modaux différents, correctement construits (sans \emph{to}, sans
\emph{-s} à la troisième personne) : 4 points. Un \emph{should} ou \emph{ought to have} +
participe, employé pour le regret : 2 points. Cohérence des temps du récit : 2 points.
Richesse et exactitude du lexique : 2 points.}

\end{document}
